\begin{table}[tbp]
\centering\small
\setlength{\tabcolsep}{4pt}
\renewcommand{\arraystretch}{1.12}
\caption{TV--persistence experiment: component comparisons.}
\label{tab:legacy_ablation}
\begin{tabular}{lrrr}
\toprule
Variant & Lower MSE & \shortstack{Median\\PRE (\%)} & \shortstack{Fused adjacencies\\(asset median)} \\
\midrule
Ridge & 0/15 & +0.00 & 0.34 \\
Nonnegative increments & 0/15 & +0.00 & 0.34 \\
TV increments & 1/15 & -0.09 & 16.47 \\
Nonnegative TV & 1/15 & -0.09 & 16.47 \\
Ridge + persistence & 9/15 & +6.80 & 0.64 \\
Nonnegative + persistence & 9/15 & +6.80 & 0.64 \\
TV + persistence & 9/15 & +6.72 & 21.56 \\
Full TV + persistence & 9/15 & +6.72 & 21.56 \\
Persistence & 4/15 & -197.57 & 29.00 \\
\bottomrule
\end{tabular}
\par\smallskip
\begin{minipage}{0.98\linewidth}\footnotesize
\textit{Notes:} Results use $H=30$, expanding estimation, and the TV experiment's own ridge reference. They are separate from the MCPB component comparisons. The lower-MSE count is out of 15 assets. The last column takes the median across assets of the mean number of fused neighboring increments, with a maximum of 29. Tuning settings remain those of the separate experiment.
\end{minipage}
\end{table}
